\documentclass{article}%
\usepackage[T1]{fontenc}%
\usepackage[utf8]{inputenc}%
\usepackage{lmodern}%
\usepackage{textcomp}%
\usepackage{lastpage}%
\usepackage{geometry}%
\geometry{margin=1in,headheight=14pt,headsep=25pt}%
\usepackage{hyperref}%
\usepackage{enumitem}%
\usepackage{fancyhdr}%
\usepackage{xcolor}%
\usepackage{url}%
\usepackage{breakurl}%
%

            \hypersetup{
                colorlinks=true,
                linkcolor=blue,
                filecolor=magenta,
                urlcolor=blue
            }
            \pagestyle{fancy}
            \fancyhf{}
            \rhead{Generated on \today}
            \cfoot{\thepage}
            
            % Configure enumeration settings
            \setlist[enumerate,1]{label=\arabic*.}
            \setlist[enumerate,2]{label=\alph*.}
            \setlist[enumerate,3]{label=\Alph*.}
            \setlist[enumerate]{itemsep=0.5em}
            
            % Define a command for red text
            \newcommand{\incorrect}[1]{\textcolor{red}{#1}}
        %
\lhead{\$L\^{}1\${-}Regression}%
\title{Practice Questions for \$L\^{}1\${-}Regression}%
\author{Generated by Scribe.AI}%
\date{\today}%
%
\begin{document}%
\normalsize%
\maketitle%
\section{Questions}%
\label{sec:Questions}%
\begin{enumerate}%
\item In $L^1$-regression, what is the primary objective of the optimization problem, and how does this objective differ from that of $L^2$-regression?%
\begin{enumerate}%
\item To minimize the sum of squared errors between the predicted and actual values, similar to $L^2$-regression.%
\item To minimize the sum of absolute errors between the predicted and actual values, unlike $L^2$-regression which minimizes squared errors.%
\item To maximize the sum of absolute errors between the predicted and actual values.%
\item To minimize the maximum absolute error between any single predicted and actual value.%
\item To find the model parameters that perfectly fit all data points, regardless of error.%
\end{enumerate}%
\vspace{1em}%
\item How does the choice of the loss function ($L^1$ vs. $L^2$) affect the sensitivity of the regression model to outliers in the dataset?%
\begin{enumerate}%
\item Both $L^1$ and $L^2$ are equally sensitive to outliers.%
\item $L^1$ is more sensitive to outliers than $L^2$.%
\item $L^2$ is more sensitive to outliers than $L^1$.%
\item $L^1$ is completely insensitive to outliers, while $L^2$ is highly sensitive.%
\item The sensitivity to outliers depends on the specific dataset and not the loss function.%
\end{enumerate}%
\vspace{1em}%
\item What is a key computational challenge associated with solving the $L^1$-regression problem compared to $L^2$-regression, and why does this challenge arise?%
\begin{enumerate}%
\item $L^1$-regression is computationally faster than $L^2$-regression due to its simpler loss function.%
\item $L^1$-regression involves a non-differentiable loss function, making it more challenging to solve using gradient-based methods.%
\item $L^1$-regression requires significantly more memory than $L^2$-regression.%
\item $L^1$-regression is inherently unstable and prone to numerical errors.%
\item There is no significant computational difference between $L^1$ and $L^2$ regression.%
\end{enumerate}%
\vspace{1em}%
\item Describe the geometric interpretation of the solution to the $L^1$-regression problem, and how does it differ from the geometric interpretation of the $L^2$-regression solution?%
\begin{enumerate}%
\item The $L^1$ solution minimizes the sum of squared distances from the data points to the regression line, while the $L^2$ solution minimizes the sum of absolute distances.%
\item The $L^2$ solution is the orthogonal projection of the data points onto the regression line, while the $L^1$ solution is not necessarily an orthogonal projection.%
\item Both $L^1$ and $L^2$ solutions are always identical, regardless of the data.%
\item The $L^1$ solution is always a point within the convex hull of the data points, while the $L^2$ solution may lie outside.%
\item There is no meaningful geometric interpretation for either $L^1$ or $L^2$ regression solutions.%
\end{enumerate}%
\vspace{1em}%
\item In the context of $L^1$-regression, explain how the choice of the loss function ($L^1$ norm) affects the sensitivity of the regression model to outliers in the dataset, and compare this to the sensitivity of an $L^2$-regression model to outliers.%
\begin{enumerate}%
\item The $L^1$ loss function is less sensitive to outliers than the $L^2$ loss function because it uses the absolute value of the errors, which gives less weight to large errors compared to the squared errors used in $L^2$ regression.%
\item The $L^1$ loss function is more sensitive to outliers than the $L^2$ loss function because it uses the absolute value of the errors, which gives more weight to large errors compared to the squared errors used in $L^2$ regression.%
\item Both $L^1$ and $L^2$ loss functions are equally sensitive to outliers because they both aim to minimize the difference between the predicted and actual values.%
\item The sensitivity to outliers depends on the specific dataset and not on the choice of the loss function ($L^1$ or $L^2$).%
\item The $L^1$ loss function is only sensitive to outliers when the number of outliers is greater than half the number of data points.%
\end{enumerate}%
\vspace{1em}%
\item In the context of $L^1$-regression, what is the primary objective of the optimization problem, and how does this objective differ from that of $L^2$-regression?%
\begin{enumerate}%
\item To minimize the sum of absolute errors, unlike $L^2$-regression which minimizes the sum of squared errors.%
\item To minimize the sum of squared errors, similar to $L^2$-regression.%
\item To maximize the sum of absolute errors.%
\item To maximize the sum of squared errors.%
\item To minimize the maximum absolute error.%
\end{enumerate}%
\vspace{1em}%
\item How does the choice of the loss function ($L^1$ vs. $L^2$) affect the sensitivity of the regression model to outliers in the dataset?%
\begin{enumerate}%
\item $L^1$-regression is more sensitive to outliers than $L^2$-regression.%
\item $L^2$-regression is more sensitive to outliers than $L^1$-regression.%
\item Both methods are equally sensitive to outliers.%
\item $L^1$-regression is less sensitive to outliers than $L^2$-regression only when the number of data points is large.%
\item The sensitivity depends on the specific dataset and not the loss function.%
\end{enumerate}%
\vspace{1em}%
\item What is a key computational challenge associated with solving the $L^1$-regression problem compared to $L^2$-regression, and why does this challenge arise?%
\begin{enumerate}%
\item The $L^1$-regression problem is always computationally faster than $L^2$-regression.%
\item The $L^1$-regression problem is non-differentiable at zero, leading to computational difficulties.%
\item The $L^2$-regression problem is non-differentiable, leading to computational difficulties.%
\item Both problems have similar computational complexities.%
\item The $L^1$-regression problem requires more memory than $L^2$-regression.%
\end{enumerate}%
\vspace{1em}%
\item Describe the geometric interpretation of the solution to the $L^1$-regression problem, and how does it differ from the geometric interpretation of the $L^2$-regression solution?%
\begin{enumerate}%
\item The $L^1$ solution minimizes the sum of distances to the fitted line, while the $L^2$ solution minimizes the sum of squared distances.%
\item Both solutions have the same geometric interpretation.%
\item The $L^1$ solution finds the median, while the $L^2$ solution finds the mean.%
\item The $L^1$ solution is always a vertex of the feasible region, while the $L^2$ solution is not necessarily a vertex.%
\item The $L^1$ solution is always unique, while the $L^2$ solution may not be unique.%
\end{enumerate}%
\vspace{1em}%
\item Given the data points (1, 2), (2, 4), (3, 5), and using the $L^1$-regression formulation from Slide 5, what is the objective function to minimize?%
\begin{enumerate}%
\item $ \min_{x_1, x_2} \sum_{i=1}^3 (y_i - (x_1 + x_2))^2$%
\item $ \min_{x_1, x_2} \sum_{i=1}^3 |y_i - (x_1 + x_2)|$%
\item $ \min_{x_1, x_2, t_1, t_2, t_3} (t_1 + t_2 + t_3)$ subject to $-t_i \le y_i - (x_1 + x_2) \le t_i$ for $i = 1, 2, 3$%
\item $ \min_{x_1, x_2} \max_{i=1,2,3} |y_i - (x_1 + x_2)|$%
\item $ \min_{x_1, x_2} \sum_{i=1}^3 (y_i - (x_1 + x_2))^2$%
\end{enumerate}%
\vspace{1em}%
\end{enumerate}

%
\newpage%
\section{Answers}%
\label{sec:Answers}%
\begin{enumerate}%
\item In $L^1$-regression, what is the primary objective of the optimization problem, and how does this objective differ from that of $L^2$-regression?%
\begin{enumerate}%
\item \incorrect{This describes the objective of $L^2$-regression, not $L^1$-regression.}%
\item The primary objective in $L^1$-regression is to minimize the sum of the absolute differences between the predicted and actual values. This contrasts with $L^2$-regression, which aims to minimize the sum of the *squares* of these differences.%
\item \incorrect{The goal is to *minimize*, not maximize, the error.}%
\item \incorrect{This describes a different type of regression problem focused on minimizing the maximum error.}%
\item \incorrect{Perfect fitting is generally not possible and not the goal of regression; the aim is to find a good balance between fitting the data and avoiding overfitting.}%
\end{enumerate}%
\vspace{1em}%
\item How does the choice of the loss function ($L^1$ vs. $L^2$) affect the sensitivity of the regression model to outliers in the dataset?%
\begin{enumerate}%
\item \incorrect{This is incorrect; the $L^1$ and $L^2$ loss functions have different sensitivities to outliers.}%
\item \incorrect{This is the opposite of the correct answer.}%
\item The $L^2$ loss function (sum of squared errors) is more sensitive to outliers because squaring large errors amplifies their influence on the model parameters.  $L^1$ (sum of absolute errors) is more robust to outliers because it treats all errors linearly.%
\item \incorrect{While $L^1$ is less sensitive, it is not completely insensitive to outliers.}%
\item \incorrect{The choice of loss function significantly impacts the sensitivity to outliers.}%
\end{enumerate}%
\vspace{1em}%
\item What is a key computational challenge associated with solving the $L^1$-regression problem compared to $L^2$-regression, and why does this challenge arise?%
\begin{enumerate}%
\item \incorrect{The opposite is true; $L^1$ regression is computationally more challenging.}%
\item The $L^1$ loss function is not differentiable at zero, making it difficult to solve using standard gradient-based optimization techniques.  $L^2$ regression has a differentiable loss function, leading to simpler solutions via normal equations or gradient descent.%
\item \incorrect{Memory requirements are not a primary distinguishing factor between the two methods.}%
\item \incorrect{Both methods can be numerically stable when implemented correctly.}%
\item \incorrect{There is a significant computational difference due to the differentiability of the loss function.}%
\end{enumerate}%
\vspace{1em}%
\item Describe the geometric interpretation of the solution to the $L^1$-regression problem, and how does it differ from the geometric interpretation of the $L^2$-regression solution?%
\begin{enumerate}%
\item \incorrect{This reverses the descriptions of $L^1$ and $L^2$ solutions.}%
\item In $L^2$-regression, the solution corresponds to the orthogonal projection of the data points onto the regression line.  In $L^1$-regression, the solution is found at the intersection of medians, and it's not necessarily an orthogonal projection.%
\item \incorrect{The solutions are generally different, except in special cases.}%
\item \incorrect{While the $L^1$ solution often lies within the convex hull, this is not a defining characteristic that distinguishes it from $L^2$.}%
\item \incorrect{Both methods have meaningful geometric interpretations.}%
\end{enumerate}%
\vspace{1em}%
\item In the context of $L^1$-regression, explain how the choice of the loss function ($L^1$ norm) affects the sensitivity of the regression model to outliers in the dataset, and compare this to the sensitivity of an $L^2$-regression model to outliers.%
\begin{enumerate}%
\item The $L^1$ norm (sum of absolute errors) is less sensitive to outliers than the $L^2$ norm (sum of squared errors).  Large errors are penalized less severely in $L^1$ regression, making it more robust to outliers that might disproportionately influence the $L^2$ regression line.%
\item \incorrect{This statement incorrectly describes the effect of the $L^1$ loss function.  The absolute value penalizes large errors less than squaring them.}%
\item \incorrect{This is incorrect. The $L^1$ and $L^2$ loss functions have different sensitivities to outliers due to the different ways they penalize errors.}%
\item \incorrect{While the dataset's characteristics play a role, the choice of loss function significantly impacts outlier sensitivity. $L^1$ is inherently more robust.}%
\item \incorrect{The sensitivity of $L^1$ regression to outliers is not dependent on a specific threshold of the number of outliers relative to the data points.}%
\end{enumerate}%
\vspace{1em}%
\item In the context of $L^1$-regression, what is the primary objective of the optimization problem, and how does this objective differ from that of $L^2$-regression?%
\begin{enumerate}%
\item The primary objective of $L^1$-regression is to minimize the sum of the absolute differences between the observed and predicted values. This contrasts with $L^2$-regression, which aims to minimize the sum of the squares of these differences.%
\item \incorrect{This describes the objective of $L^2$-regression, not $L^1$-regression.}%
\item \incorrect{The goal is to minimize, not maximize, the error.}%
\item \incorrect{This describes the objective of $L^2$-regression, not $L^1$-regression.}%
\item \incorrect{This describes a different type of regression problem (minimax).}%
\end{enumerate}%
\vspace{1em}%
\item How does the choice of the loss function ($L^1$ vs. $L^2$) affect the sensitivity of the regression model to outliers in the dataset?%
\begin{enumerate}%
\item \incorrect{This is incorrect; $L^1$ regression is less sensitive to outliers.}%
\item Because the $L^2$ loss function squares the errors, outliers have a disproportionately large influence on the model's parameters. The $L^1$ loss function, using absolute errors, is less sensitive to outliers.%
\item \incorrect{This is incorrect; $L^1$ regression is less sensitive to outliers.}%
\item \incorrect{The sensitivity difference holds regardless of dataset size.}%
\item \incorrect{The choice of loss function significantly impacts outlier sensitivity.}%
\end{enumerate}%
\vspace{1em}%
\item What is a key computational challenge associated with solving the $L^1$-regression problem compared to $L^2$-regression, and why does this challenge arise?%
\begin{enumerate}%
\item \incorrect{This is incorrect; $L^1$ regression is generally more computationally expensive.}%
\item The absolute value function in the $L^1$ loss is not differentiable at zero, making it more challenging to find the minimum using gradient-based methods compared to the smooth $L^2$ loss function.%
\item \incorrect{This is incorrect; the $L^2$ loss function is differentiable.}%
\item \incorrect{This is incorrect; $L^1$ regression is computationally more challenging.}%
\item \incorrect{Memory requirements are not the primary computational difference.}%
\end{enumerate}%
\vspace{1em}%
\item Describe the geometric interpretation of the solution to the $L^1$-regression problem, and how does it differ from the geometric interpretation of the $L^2$-regression solution?%
\begin{enumerate}%
\item In $L^1$ regression, the solution minimizes the sum of the vertical distances from the data points to the regression line. In $L^2$ regression, it minimizes the sum of the squared vertical distances.%
\item \incorrect{The geometric interpretations are different.}%
\item \incorrect{While this is true in simple cases, it's not the complete geometric interpretation.}%
\item \incorrect{While often true, it's not always the case for $L^1$ solutions.}%
\item \incorrect{Neither solution is guaranteed to be unique in all cases.}%
\end{enumerate}%
\vspace{1em}%
\item Given the data points (1, 2), (2, 4), (3, 5), and using the $L^1$-regression formulation from Slide 5, what is the objective function to minimize?%
\begin{enumerate}%
\item \incorrect{This minimizes the sum of squared errors, which is the $L^2$-regression objective, not $L^1$.}%
\item \incorrect{While this represents the sum of absolute errors, it's not in a form suitable for linear programming. The linear programming formulation requires the introduction of slack variables.}%
\item The $L^1$-regression aims to minimize the sum of absolute errors.  Slide 5 shows this is achieved by introducing slack variables $t_i$ and minimizing their sum, subject to constraints ensuring each $t_i$ bounds the absolute error.%
\item \incorrect{This minimizes the maximum absolute error, which is the $L^\infty$-regression objective.}%
\item \incorrect{This is the same as option A, minimizing the sum of squared errors ($L^2$-regression).}%
\end{enumerate}%
\vspace{1em}%
\end{enumerate}

%
\end{document}